\subsection{与滤过模相伴的分次模}

设 G 为交换群（用加法记号），\((G_{n})\) 为 G 上的一个滤过。记

\[
\begin{array}{l} \operatorname{gr} _ {n} (\mathbf {G}) = \mathbf {G} _ {n} / \mathbf {G} _ {n + 1} \quad \text { for } \quad n \in \mathbf {Z} \\ \operatorname{gr} (\mathbf {G}) = \bigoplus_ {n \in \mathbf {Z}} \operatorname{gr} _ {n} (\mathbf {G}). \end{array}\tag{8}
\]

	于是，交换群 \(\operatorname{gr}(\mathbf{G})\) 是一个 \(\mathbf{Z}\) 型分次群，称为与滤过群 \(\mathbf{G}\) 相伴的分次群；其 \(n\) 次齐次元属于 \(\operatorname{gr}(\mathbf{G})\)，并且是 \(\operatorname{gr}_n(\mathbf{G})\) 的元素。

现在设 A 为滤过环，\((\mathbf{A}_{n})\) 为其滤过，E 为滤过 A-模，\((\mathbf{E}_{n})\) 为其滤过。对所有 \(p \in Z\)、\(q \in Z\)，定义映射

\[
\operatorname{gr} _ {p} (\mathrm{A}) \times \operatorname{gr} _ {q} (\mathrm{E}) \rightarrow \operatorname{gr} _ {p + q} (\mathrm{E})\tag{9}
\]

	定义如下：给定 \(\alpha \in \mathrm{gr}_p(\mathbf{A})\)、\(\xi \in \mathrm{gr}_q(\mathbf{E})\)，取两个代表元 \(a, a'\) 代表 \(\alpha\)，以及两个代表元 \(x, x'\) 代表 \(\xi\)，则 \(ax \in \mathrm{E}_{p + q}\)、\(a'x' \in \mathrm{E}_{p + q}\)，且 \(ax \equiv a'x'\) (模 \(\mathrm{E}_{p + q + 1}\))，因为

\[
a x - a ^ {\prime} x ^ {\prime} = (a - a ^ {\prime}) x + a ^ {\prime} (x - x ^ {\prime})
\]

	又有 \(a - a' \in \mathbf{A}_{p+1}\) 和 \(x - x' \in \mathbf{E}_{q+1}\)，故上述断言由第 1 节的公式 (3) 得出。因此，我们可以把 \(\alpha\xi\) 记作

\[
\mathbf {E} _ {p + 1} / \mathbf {E} _ {p + q + 1} = \operatorname{gr} _ {p + q} (\mathbf {E})
\]

	任取代表元 \(a \in \alpha\) 和任取代表元 \(x \in \xi\) 的乘积 ax 在上述商中的类。显然，映射 (9) 是 Z-双线性的；由线性性得到一个 Z-双线性映射

\[
\operatorname{gr} (\mathrm{A}) \times \operatorname{gr} (\mathrm{E}) \rightarrow \operatorname{gr} (\mathrm{E}).\tag{10}
\]

	若先将这个定义应用于 \(E = A_{s}\) 的情形，则映射 (10) 是 \(\mathrm{gr}(A)\) 上的内部合成律；直接验证可知它具有结合性，并且有一个单位元，即 A 的单位元在 \(\mathrm{gr}_{0}(A)\) 中的典范像。因此，它在 \(\mathrm{gr}(A)\) 上定义了一个环结构，而分次 \((\mathrm{gr}_{n}(A))_{n\in\mathbf{Z}}\) 按定义与此结构相容。这样定义的（Z 型）分次环 \(\mathrm{gr}(A)\) 称为与滤过环 A 相伴的分次环；若 A 交换，它显然也交换；\(\mathrm{gr}_{0}(A)\) 是 \(\mathrm{gr}(A)\) 的子环。另一方面，映射 (10) 是 \(\mathrm{gr}(A)\)-模 \(\mathrm{gr}(E)\) 上的外部合成律，模公理显然满足，而分次 \((\mathrm{gr}_{n}(E))_{n\in\mathbf{Z}}\) 在 \(\mathrm{gr}(E)\) 上显然与这个模结构相容。这样定义的分次 \(\mathrm{gr}(A)\)-模 \(\mathrm{gr}(E)\)（Z 型）称为与滤过 A-模 E 相伴的分次模。

\textbf{例}\par

\begin{enumerate}
\def\labelenumi{(\arabic{enumi})}
\item
	  设 A 为交换环，t 为 A 中的一个非零因子。给 A 赋予 \((t)\)-进滤过（第 1 节，例 3）。则与之相伴的分次环 \(\mathrm{gr}(\mathbf{A})\) 典范同构于多项式环 \((\mathbf{A}/(t))[\mathbf{X}]\)。因为 \(\mathrm{gr}_{n}(\mathbf{A}) = 0\) 对 n \textless{} 0 成立，且按定义环 \(\mathrm{gr}_{0}(\mathbf{A})\) 就是环 \(\mathbf{A}/(t)\)。现在注意到，由于关于 t 的假设，关系 \(at^{n} \equiv 0 \pmod{t^{n+1}}\) 等价于 \(a \equiv 0 \pmod{t}\)；若 \(\tau\) 是 t 在 \(\mathrm{gr}_{1}(\mathbf{A})\) 中的典范像，则 \(\mathrm{gr}_{n}(\mathbf{A})\) 的每个元素都可唯一写成 \(\alpha\tau^{n}\) 的形式，其中 \(\alpha \in \mathrm{gr}_{0}(\mathbf{A})\)；断言得证。
\item
  设 \(\mathbf{K}\) 为交换环，A 为形式幂级数环

\[
\mathrm{K} \left[ \left[ \mathrm{X} _ {1}, \dots , \mathrm{X} _ {r} \right] \right]
\]

(《代数》，第 IV 章，§ 5)，m 为 A 的理想，其元素是没有常数项的形式幂级数。给 A 赋予 m-进滤过（第 1 节，例 3）；若 \(M_{1}, \ldots, M_{s}\) 是 \(X_{1}, \ldots, X_{r}\) 中总次数为 n - 1 的互异单项式，则显然每个总阶为 \(\omega(u) \geqslant n\) 的形式幂级数 u

	(同上，第 2 节)，可写成 \(\sum_{k=1} u_k M_k\)，其中 \(u_k\) 属于 \(m\)；可以看出，\(m^n\) 是这样一个形式幂级数 \(u\) 的集合：它满足 \(\omega(u) \geqslant n\)，这说明 \(\omega\) 是 m-进滤过的阶函数。于是显然，对每个形式幂级数 \(u \in m^n\)，存在唯一一个次数为 \(n\) 的齐次多项式（关于 \(X_i\)），它与 \(u\) 模 \(m^{n+1}\) 同余，即 \(n\) 次项在 \(u\) 中的和；因此 \(\text{gr}(A)\) 典范同构于多项式环 \(K[X_1, \ldots, X_f]\)。

\item
	  更一般地，设 A 为交换环，b 为 A 的理想，并给 A 赋予 b-进滤过。若记 \(\mathbf{B} = \mathbf{gr}_{0}(\mathbf{A})\)、\(\mathbf{F} = \mathbf{gr}_{1}(\mathbf{A}) = \mathbf{b}/\mathbf{b}^{2}\)，我们知道（《代数》，第 III 章），B-模 F 到自身的恒等映射可以唯一延拓为同态 u，使其把对称代数 \(\mathbf{S}(\mathbf{F})\)（其中 F 是 B-模）映到 B-代数 \(\mathbf{gr}(\mathbf{A})\)；由 \(\mathbf{gr}(\mathbf{A})\) 的定义可知，u 是分次代数的满同态；当 \(n \geqslant 1\) 时，\(\mathbf{gr}_{n}(\mathbf{A})\) 的每个元素都是元素模 \(b^{n+1}\) 的类之和，这些元素形如 \(y = x_{1}x_{2}\ldots x_{n}\)，其中 \(x_{i} \in b (1 \leqslant i \leqslant n)\)；若 \(\xi_{i}\) 是 \(x_{i} \mod b^{2}\) 的类，则显然 y 模 \(b^{n+1}\) 的类就是元素 \(u(\xi_{1})\ldots u(\xi_{n})\)，断言得证。特别地，B-模 F 的每个生成系都是 B-代数 \(\mathbf{gr}(\mathbf{A})\) 的生成系。
\end{enumerate}

现在若 E 是 A-模且给 E 赋予 b-进滤过，同样可见，分次 gr(A)-模 gr(E) 由 gr₀(E) = E/bE 生成。确切地说，将 gr(A)-模 gr(E) 上的外部合成律限制到 gr(A) × gr₀(E) 后，得到一个从 gr(A) × gr₀(E) 到 gr(E) 的 Z-双线性映射 φ；此外，gr(A) 是一个 (gr₀(A), gr₀(A))-双模，而 gr₀(E) 是一个 gr₀(A)-模；直接验证可知，对 α ∈ gr(A)、α₀ ∈ gr₀(A)、ξ ∈ gr₀(E)，

\[
\phi (\alpha \alpha_ {0}, \xi) = \phi (\alpha , \alpha_ {0} \xi)
\]

从而 \(\phi\) 定义了一个满的 \(\mathbf{gr}_0(\mathbf{A})\)-线性映射

\[
\gamma_ {\mathrm{E}} \colon \operatorname{gr} (\mathrm{A}) \otimes_ {\mathbf {g r} _ {0} (\mathrm{A})} \operatorname{gr} _ {0} (\mathrm{E}) \rightarrow \operatorname{gr} (\mathrm{E})\tag{11}
\]

称为典范映射。

	* (4) 设 K 为交换环，g 为 K 上的李代数，U 为 g 的包络代数。在 U 上定义递增滤过 \((\mathrm{U}_{n})_{n\in\mathbf{Z}}\)：取 \(U_{n}=\{0\}\)（当 n\textless0 时），并令 \(U_{n}\)（当 \(n\geqslant0\) 时）为 U 中可表示为至多 n 个 g 中元素之积之和的元素所成的集合；于是

	\(U_{0} = K\)，且 \(\text{gr}(U)\) 是交换 K-代数（《李群与李代数》，第 I 章，§ 2，第 6 节）。从 \(g\) 到 \(\text{gr}_{1}(U) = U_{1}/U_{0}\) 的典范映射可以唯一延拓为同态 \(h\)，它把对称代数 \(S(g)\)（其中 \(g\) 是 K-模）映到 K-代数 \(\text{gr}(U)\)；同态 \(h\) 是满射；若 K-模 \(g\) 是自由的，则 \(h\) 是双射（同上，第 7 节，定理 1）。

\begin{enumerate}
\def\labelenumi{(\arabic{enumi})}
\setcounter{enumi}{4}
\tightlist
\item
  设 A 为 Z 型分次环，E 为 Z 型分次 A-模；令 \(A_{(i)}\)（分别为 \(E_{(i)}\)）表示 A（分别 E）的 i 次齐次元所成的子群。给 A 和 E 赋予与其分次相伴的滤过（第 1 节，例 1），并令 \(A'\) 和 \(E'\) 表示由此得到的滤过环和滤过 \(A'\)-模。于是显然，Z-线性映射 \(A \to \text{gr}(A')\) 把 \(A_{(n)}\) 中的元素映到其在

\[
\operatorname{gr} _ {n} (\mathbf {A}) = \left(\bigoplus_ {i \geqslant n} \mathbf {A} _ {(i)}\right) / \left(\bigoplus_ {i \geqslant n + 1} \mathbf {A} _ {(i)}\right)
\]

是分次环同构。类似地定义一个典范分次 A-模同构 \(\mathbf{E} \to \mathrm{gr}(\mathbf{E}')\)。
\end{enumerate}

	命题 1. 设 A 为滤过环，\((\mathbf{A}_n)_{n\in \mathbf{Z}}\) 为其滤过，\(v\) 为其阶函数。假设 \(\operatorname {gr}(\mathbf{A})\) 是无零因子的环。那么，对于任意有序对元素 \(a,b\)，在环 \(\mathbf{B} = \bigcup_{n\in \mathbf{Z}}\mathbf{A}_n,v(ab) = v(a) + v(b)\) 中，所述等式成立。

	由于 \(n = \bigcap_{n \in \mathbf{Z}} A_n\) 是环 \(B\) 的双边理想，当 \(v(a)\) 或 \(v(b)\) 等于 \(+\infty\) 时上述公式成立。否则，\(v(a) = r\) 和 \(v(b) = s\) 都是整数；按定义，\(\alpha\) 是 \(a \bmod A_{r+1}\) 的类，\(\beta\) 是 \(b \bmod A_{s+1}\) 的类，且二者都 \(\neq 0\)，因此由假设，\(\alpha \beta \neq 0\) 于 \(\operatorname{gr}(A)\) 中，从而 \(ab \notin A_{r+s+1}\)；而 \(ab \in A_{r+s}\)，

\[
v (a b) = v (a) + v (b).
\]

推论。设 A 为滤过环，\((\mathbf{A}_{n})_{n\in\mathbf{Z}}\) 为其滤过；令 \(B=\bigcup_{n\in Z}A_{n}\)、\(n=\bigcap_{n\in Z}A_{n}\)。若环 gr(A) 无零因子，则环 B/n 也无零因子。

	若 \(a\) 和 \(b\) 是 \(B\) 中不属于 \(\mathfrak{n}\) 的元素，则 \(v(a) \neq +\infty\) 且 \(v(b) \neq +\infty\)，从而 \(v(ab) \neq +\infty\)，因此 \(ab \notin \mathfrak{n}\)。

注意，环 A 可以是整环，且滤过 \((A_{n})\) 穷竭并分离，但 gr(A) 未必是整环（习题 2）。

	注。设 G 为不必交换的群，带有滤过 \((\mathrm{G}_{n})_{n\in\mathbf{Z}}\)，使得 \(G_{n+1}\) 是 \(G_{n}\) 的正规子群，对所有 \(n\in Z\)；仍令 \(\mathrm{gr}_{n}(\mathrm{G})=\mathrm{G}_{n}/\mathrm{G}_{n+1}\)。族 \((\mathrm{gr}_{n}(\mathrm{G}))_{n\in\mathbf{Z}}\) 的限制直积，即乘积 \(\prod_{n\in\mathbf{Z}}\mathrm{gr}_{n}(\mathrm{G})\) 中由元素 \((\xi_{n})\) 构成的子群，其中除至多有限个分量外，其余所有分量都等于单位元，也称为与 G 相伴的分次群，记为 \(\mathrm{gr}(\mathrm{G})\)。

